\chapter{Percepção do consumidor sobre \textit{Brand Equity} com base na sustentabilidade ambiental – Desenvolvimento e validação de uma escala}


\begin{flushright}
Luiz Diego Vidal Santos  \footnote[1]{Programa de Pós-graduação em Ciências da Propriedade Intelectual-PPGPI/Universidade Federal de Sergipe-UFS. *E-mail:  \href{mailto:vidal.center@academico.ufs.br}{vidal.center@academico.ufs.br}}
\\[0.2cm] 
Francisco Sandro Rodrigues Holanda \footnote[2]{Departamento de Engenharia Agronômica/Universidade Federal de Sergipe, São Cristóvão, Brasil. E-mail \href{mailto:fholanda@academico.ufs.br}{fholanda@academico.ufs.br}}
\\[0.2cm] 
Catuxe Varjão de Santana Oliveira \footnote[3]{Departamento de Ciências da Computação/Instituto Federal de Sergipe, Lagarto, Brasil. E-mail \href{mailto:catuxe.oliveira@ifs.edu.br}{catuxe.oliveira@ifs.edu.br}}
\\[0.2cm] 
Pedro Vinícius Bertulino de Menezes \footnote[2]{Departamento de Engenharia Agronômica/Universidade Federal de Sergipe, São Cristóvão, Brasil. E-mail \href{mailto:pedrovinicius18@academico.ufs.br}{pedrovinicius18@academico.ufs.br}}
\\[0.2cm] 
Ana Paula Schervinski Villwock \footnote[2]{Departamento de Engenharia Agronômica/Universidade Federal de Sergipe, São Cristóvão, Brasil. E-mail \href{mailto:ana.agronomia@gmail.com}{ana.agronomia@gmail.com}}
\\[0.2cm] 
Alceu Pedrotti \footnote[2]{Departamento de Engenharia Agronômica/Universidade Federal de Sergipe, São Cristóvão, Brasil. E-mail \href{mailto:alceupedrotti@gmail.com}{alceupedrotti@gmail.com}}
\end{flushright}

% resumo em português
\begin{resumoumacoluna}
As marcas gradualmente se tornaram a dimensão central e o ativo estratégico do \textit{Branding} para organizações de todos os tamanhos e hoje muitas empresas adotam várias formas de atividades de marketing verde como parte de suas estratégias. O objetivo deste estudo foi desenvolver e validar as propriedades psicométricas de precisão de uma escala para avaliação do valor da marca, com base na sustentabilidade ambiental. A escala foi validada adotando-se os padrões descritos na Norma American Psychological Association (APA). Oito empresas foram selecionadas dentre as que desenvolvem seus negócios sustentáveis como amostras para testes de validação. Os questionários foram aplicados a 263 potenciais consumidores. Foi construída e validada uma escala com 2 itens, dividida em quatro dimensões: Percepção da qualidade da marca BQP (6 questões); SBP de posicionamento estratégico de marca (10 questões); Disponibilidade para comprar PI (3 questões) e Retail Innovation RI (3 questões). Os itens do questionário apresentaram consistência interna aceitável quanto aos seus construtos. Matrix é interpretado em 4 dimensões. A análise fatorial exploratória assumiu poucos carregamentos cruzados, corroborando os pressupostos iniciais dos instrumentos adaptados. O instrumento obteve bons índices de ajuste interno e medidas externas. Finalmente, a análise da teoria da resposta ao item (TRI) não mostrou nenhum efeito piso ou teto; no entanto, para os dados do DIF, a escala invariavelmente funciona para pessoas mais velhas e mais jovens para quase todos os itens, exceto o item RI 04. Este artigo preenche a lacuna na literatura ao desenvolver uma escala para avaliar a percepção do consumidor sobre o valor da marca e a sustentabilidade ambiental, melhorando a compreensão da experiência do cliente com o valor da marca, explorando várias dimensões da percepção do consumidor sobre marcas sustentáveis.
\vspace{\onelineskip}

\noindent
\textbf{Palavras-chave}: Sustentabilidade florestal; Sustentável; Validação de escala; Instrumento psicométrico; Floresta Amazônica.

\noindent
Artigo sob avaliação no periódico: Journal of Marketing Analytics
%\href{https://doi.org/10.53660/CLM-106-127B}{\color{blue}{\underline{https://doi.org/10.53660/CLM-106-127B}}}

\noindent
\begin{figure}[H]
  \centering
  \caption*{Resumo gráfico 2º capítulo}
  \includegraphics[scale=0.15]{Imagens/resumo_2.pdf}
  \fonte{\textbf{Elaborado pelos autores} (2022)}
  \label{figres:03}
\end{figure}

\href{https://empreendasustentavel.com/resumo_2.html}{\color{blue}{\underline{Download: Baixe a imagem em tamanho real}}}

\end{resumoumacoluna}


\section{Introdução}

O valor da marca de empresas e o enfoque de sustentabilidade das ações empreendedoras têm sido constantemente discutidos no cenário acadêmico \cite{Moritz2018}. Esses dois fundamentos devem ser vistos como simbióticos ao planejar a mudança da imagem do negócio com foco na inovação no \textit{setlist} atual. Outros estudos reforçam que a sustentabilidade é um fator antecedente e básico para o \textit{Brand Equity} da marca \cite{Gonzalez-Lafaysse2016FacebookChain,Arrigo2017281,Ray2020AntecedentsSustainability}. Pesquisas evidenciam uma ligação entre essas variáveis que podem contribuir para as vantagens econômicas e sustentáveis de uma organização \cite{Marin-Garcia2020HowRetailing}. 

Dentro deste contexto, é fundamental compreender como os consumidores percebem e valorizam as marcas que incorporam práticas sustentáveis em suas estratégias de negócio \cite{Mehdikhani2022TheFactors}. Nesse sentido, é necessário estabelecer uma conexão entre a percepção do consumidor sobre marcas sustentáveis e o impacto dessas percepções no \textit{Brand Equity} \cite{Emari2012TheEquity}.

A importância da sustentabilidade para o \textit{Brand Equity} pode ser melhor entendida ao analisar o papel das marcas como principal ativo para o nascimento da percepção do consumidor. As marcas podem ser entendidas como rótulos, símbolos e descrições que envolvem significados e associações, resultando em um contrato flexível com o cliente na forma de prestação de um produto ou serviço, signos distintos de produtos, empresas ou serviços \cite{Woodrow2007}. Com isso, tornaram-se um importante ativo econômico para seu portador, determinando o quanto o consumidor aceita o produto da empresa signatária da marca, impactando significativamente seu desempenho comercial \cite{Singh2020CorporateAlliances}. Diante disso, diversas empresas surgem ou, quando consolidadas, adotam políticas sustentáveis em todas as esferas de contato com o consumidor.

Para estabelecer uma conexão eficiente entre a sustentabilidade e o \textit{Brand Equity}, é necessário compreender os fatores que influenciam a percepção do consumidor em relação às marcas sustentáveis. Entre esses fatores, destacam-se: a qualidade percebida, o posicionamento da marca, a intenção do consumidor para comprar e a inovação sustentável da marca \cite{Aaker1997DimensionsPersonality, Figueiredo2019PassengersBrazil}. 

A análise desses fatores permite identificar as estratégias mais eficientes para desenvolver marcas sustentáveis que atendam às expectativas do consumidor e contribuam para o sucesso comercial das empresas. Além disso, a compreensão desses fatores e sua influência na percepção do consumidor pode ser útil para o desenvolvimento de políticas e iniciativas que promovam a adoção de práticas sustentáveis por parte das empresas e estimulem o consumo consciente por parte dos consumidores.

Com base no exposto, é possível concluir que a sustentabilidade e o \textit{Brand Equity} estão intrinsecamente relacionados e devem ser considerados integradamente no planejamento e execução das estratégias empresariais. A pesquisa e o desenvolvimento de instrumentos psicométricos que avaliem a percepção do consumidor sobre marcas sustentáveis e o impacto dessas percepções no \textit{Brand Equity} são essenciais para orientar as empresas na construção de marcas fortes e sustentáveis, capazes de atender às demandas e expectativas dos consumidores e contribuir para o sucesso comercial e a sustentabilidade ambiental.

Ao considerar a necessidade de desenvolver instrumentos psicométricos eficazes para avaliar a percepção do consumidor sobre marcas sustentáveis e o impacto dessas percepções no \textit{Brand Equity}, é importante destacar a relevância da adaptação transcultural desses instrumentos. Uma vez que diferentes populações possuem diferentes valores sociais e culturais, é essencial garantir que os instrumentos desenvolvidos sejam culturalmente adequados e capazes de capturar as nuances específicas de cada contexto \cite{Zajkowski2019DifferencesProspect}.

Neste sentido, é fundamental considerar a adaptação transcultural dos instrumentos psicométricos no processo de desenvolvimento, garantindo que os conceitos e itens sejam compreendidos de maneira semelhante pelos respondentes de diferentes culturas. Para alcançar isso, é necessário realizar uma tradução e adaptação cuidadosa dos itens, considerando as características linguísticas e culturais da população-alvo \cite{Hustvedt2010EffectsApparel}.

Além disso, é importante ressaltar a importância de validar e refinar os instrumentos propostos por meio de metodologias robustas e rigorosas. A validação do instrumento deve seguir os padrões estabelecidos pela APA (American Psychological Association) e outros órgãos reguladores, garantindo a precisão e validade das propriedades psicométricas da escala desenvolvida \cite{Pandey2014E-lifestylesValidation,Boateng2018BestPrimer}.

Com o desenvolvimento e a validação de instrumentos psicométricos eficazes para avaliar a percepção do consumidor sobre marcas sustentáveis e o impacto dessas percepções no \textit{Brand Equity}, as empresas podem obter insights valiosos para orientar suas estratégias de \textit{Branding} e sustentabilidade \cite{Kolk2014SustainableDressing}. Esses insights podem ajudar as empresas a identificar as áreas em que suas práticas sustentáveis são mais valorizadas pelos consumidores, bem como aquelas em que possa haver oportunidades para melhorias e inovações.

Dessa forma, a integração da sustentabilidade e do \textit{Brand Equity} no planejamento e execução das estratégias empresariais não só contribui para o sucesso comercial das empresas, mas também para a promoção de um consumo mais consciente e responsável por parte dos consumidores. Ao adotar práticas sustentáveis e comunicá-las eficazmente aos consumidores, as empresas têm a oportunidade de se destacar no mercado, criar valor para seus stakeholders e contribuir para a construção de um futuro mais sustentável e equitativo \cite{Chernev2015DoingResponsibility}.

Em suma, a pesquisa e o desenvolvimento de instrumentos psicométricos que avaliam a percepção do consumidor sobre marcas sustentáveis e o impacto dessas percepções no \textit{Brand Equity} são fundamentais para orientar as empresas na construção de marcas fortes e sustentáveis. A adaptação transcultural e a validação rigorosa desses instrumentos são aspectos cruciais para garantir sua eficácia e aplicabilidade em diferentes contextos culturais e sociais \cite{Hur2014HowReputation}. Com essas ferramentas à intenção, as empresas podem tomar decisões mais informadas e eficazes para promover a sustentabilidade e o sucesso comercial em um mundo cada vez mais globalizado e consciente dos desafios ambientais \cite{Fatma2016MeasuringResponses}.

Buscando contemplar esta lacuna de pesquisa, o objetivo deste estudo foi investigar e desenvolver as propriedades psicométricas de precisão e validade de uma escala para avaliação de valor de marca baseada na sustentabilidade ambiental da floresta amazônica brasileira. As fases de validação seguiram a recomendação de validação para instrumentos psicométricos padrão - APA. Os padrões têm importantes implicações gerais para estudos psicométricos comparativos que investigam as percepções do consumidor. As ferramentas de avaliação de escala padronizada devem demonstrar uma interpretação semelhante dos itens do teste transculturalmente.

Esta pesquisa está organizada da seguinte forma: primeiramente, apresentamos as dimensões analisadas para o desenvolvimento da escala. Em seguida, apresentamos a metodologia adotada para validar e refinar a escala proposta. Por fim, discutimos os resultados obtidos e as considerações finais.


\section{Materiais e métodos}

\subsection{Operacionalização de Variáveis}

O valor da marca compreende uma série de dimensões que a marca pode desenvolver, gerenciar e controlar. Diferentes estruturas conceituais incluem diferentes dimensões. Então, esta pesquisa é considerada um referencial teórico para as dimensões postuladas por conscientização, lealdade, qualidade percebida e associações. As dimensões iniciais e os achados de uma \textit{Scoping Review} foram desenvolvidos por meio de tradução transcultural (\textit{'forward'}) de itens selecionados dos instrumentos originais. \cite{Aaker1997DimensionsPersonality}

Dois tradutores, um linguista e outro da área de pesquisa do experimento, foram selecionados, pois na tradução de uma escala, devem-se buscar diversas equivalências sobre o original, como cultural, semântica, técnica, conteúdo e critérios \cite{Guillemin1993Cross-culturalGuidelines}. Além disso, a literatura atual na área destaca a necessidade de se evitar a tradução literal dos itens, fato que pode ser evitado se forem atendidos apenas por meio da variabilidade dos tradutores \cite{Callegaro2018}.

\subsection{Dimensões e sua análise}

A primeira dimensão é composta por 10 (dez) questões transliteradas da escala desenvolvida por \citeonline{Grigorescu2019KeySustainability}, que analisou a percepção do entrevistado sobre o posicionamento estratégico da marca (\textit{Strategic brand positioning} - SBP) na região amazônica do Brasil. Adicionalmente, foi considerado o nível de percepção dos participantes quanto à importância da relação empresa ou \textit{startup} para a sustentabilidade ambiental na região. Foram utilizados 10 (dez) com pequenas modificações de linguagem dos 21 (vinte e um) itens da escala original.

A segunda dimensão apresenta cinco itens e compreende a percepção da qualidade da marca relacionada à sustentabilidade ambiental proposta inicialmente por cada marca (\textit{perception of brand quality} - BQP). Esse construto foi avaliado por meio de cinco fotos e itens da marca, cinco fotos da loja (digital), quatro fotos e itens do produto \textit{premium} (com valor acima da média de mercado) e cinco fotos representando a intenção do negócio em relação à sustentabilidade ambiental da floresta Amazônica. (meios de produção/extração). Das mãos das figuras/produtos que buscam contemplar essa dimensão, 5 dos 11 itens foram adaptados, tratando da percepção de qualidade da marca dos questionários desenvolvidos por \citeonline{Netemeyer_Krishnan} e \citeonline{Yoo2017AnEquity:}.

A terceira dimensão contém três itens, que buscam entender quantos participantes estariam dispostos a comprar produtos de empresas relacionadas à sustentabilidade ambiental por meio dos dados apresentados (\textit{Purchase Intention} - PI). Para mensurar a intenção de compra, os 3 (três) itens estão dispostos de acordo com \citeonline{Arnett-meiers}.

Por fim, a quarta dimensão avalia o quanto o negócio apresenta inovação no varejo (RI). A dimensão foi montada solicitando aos participantes que respondessem a quatro questões que abordavam o tema inovação, singularidade e design de marca. As questões foram adaptadas dos questionários propostos por \citeonline{Netemeyer_Krishnan} e \citeonline{Lin2015ConceptualizingTaiwan}. Desses trabalhos, foi realizada apenas a tradução dos itens.

Buscando atender aos princípios éticos em pesquisa com seres humanos, este estudo foi submetido ao Comitê de Ética em Pesquisa, recebendo o código CAAE: 55999722.1.0000.5546. Os participantes que concordaram em participar do estudo foram informados sobre os objetivos da pesquisa, coleta de dados e informações referentes ao Termo de Consentimento Livre e Esclarecido (TCLE).

\subsection{Tradução transcultural dos itens para o instrumento}

A tradução e adaptação dos itens utilizados dos questionários originais seguiram cinco etapas essenciais recomendadas pela literatura: (1) tradução do instrumento do idioma de origem para o idioma de destino, (2) análise dos protótipos e análise da versão sintetizada por juízes especialistas, (3) retrotradução para o idioma de origem e (4) estudo piloto \cite{Hambleton1993TranslatingStudies,Gjersing2010Cross-culturalConsiderations}. Para conveniência desta pesquisa, outra tarefa (fase 2) foi realizada durante a validação de conteúdo pelo CVC para a escala completa, e as fases (3) e (4) foram realizadas após a finalização da validação de conteúdo.

\subsection{Público piloto e aplicação inicial da ferramenta}

Para abranger mais regiões do país, o instrumento foi aplicado tanto presencialmente quanto virtualmente. Segundo \citeonline{Walter2016}, devido às suas vantagens, como velocidade e capacidade, a aplicação de questionários digitais tem disparado para atingir populações específicas, reduzindo os possíveis custos de deslocamento. A seleção dos entrevistados participantes foi realizada de forma aleatória para significar a variabilidade populacional no Brasil.

\subsection{Validação da escala}

Após a formulação e estruturação dos itens a serem utilizados na escala de coleta de dados, ela foi submetida à validação com base no modelo tripartido descrito nas Normas para Testes Educacionais e Psicológicos, \cite{Association2018}. As Normas apresentam cinco validações de processos para aceitar o instrumento como confiável e replicável:


\begin{enumerate}[label=(\roman*)]
- Validação de conteúdo;
- Validade baseada na estrutura interna;
- Validade baseada em relacionamentos com medidas externas;
- Validade baseada no padrão de resposta aos itens.
\end{enumerate}


\subsubsection*{\textit{Validação do conteúdo}}
O coeficiente de validade de conteúdo (CVC) foi aplicado no instrumento proposto para medir a eficácia da escala com base em avaliações qualitativas de juízes especialistas na área. Essa técnica envolve buscar o consenso de uma comunidade de profissionais que irá atribuir valores e opiniões sobre as traduções e adaptações feitas no instrumento para poder ser aperfeiçoado e contemplar os objetivos do instrumento. A avaliação da escala pelos juízes foi realizada às cegas; portanto, a confidencialidade dos participantes garante maior fidelidade aos resultados ao final do experimento, permite melhor exposição dos resultados e mantém seu entendimento \cite{Hasson2000}. Após validação dos itens, a randomização foi realizada no software Excel, e a coleta de dados foi realizada de março a maio de 2022.

\subsubsection*{\textit{Validade baseada na estrutura interna}}

A validade baseada na estrutura interna reflete o grau onde a estrutura de correlações entre os itens corresponde ao que o instrumento planeja medir \cite{Association2018}. Portanto, nesta pesquisa, foi realizado um AFE para avaliar a estrutura nomeada de fatores da escala proposta. A análise foi implementada usando uma matriz policrômica e o método de extração utilizou mínimos quadrados ponderados na diagonal total (RDWLS) \cite{Asparouhov2010}. A decisão sobre o número de fatores a serem retidos foi realizada por meio da técnica de análise paralela com permutação aleatória dos dados observados \cite{Timmerman2011}.

Para determinar o número de fatores de retenção e facilitar a interpretação dos fatores \cite{Damasio2012}, foi adotado um método de rotação semelhante ao  Robust Promin \cite{Ferrando2018AssessingAnalysis}. Buscando um melhor refinamento da escala, assumiu-se que, para um item ser aceito, ele deve ter carga fatorial acima de 0.60 e cargas cruzadas entre itens menores que 0.20 \cite{Howard2016AImprove}.


Buscando avaliar a qualidade da adequação do modelo e avaliar o quão bem ele pode reproduzir a estrutura de covariância ou a estrutura correlacional das variáveis, foram utilizadas as seguintes estratégias: a raiz do erro quadrático médio de aproximação (RMSEA), o índice de Tucker-Lewis (TLI) e índice de ajuste comparativo (CFI). Segundo \citeonline{Brown2015}, o valor RMSEA deve ser menor que 0.08 e o intervalo de confiança não pode chegar a 0.10. Os valores de CFI e TLI devem ser maiores que 0.90. Preferencialmente 0.95. A confiabilidade do modelo, ou seja, quão consistente e reprodutível ele é na mensuração das variáveis latentes pretendidas, foi testada por meio do indicador de confiabilidade composta, que deve ter valores maiores ou iguais a 0.70 \cite{Hair2009}.

A estabilidade do fator foi avaliada pelo índice H \cite{Lorenzo-Seva2019}. O índice H propõe avaliar quão bem um conjunto de itens representa um fator comum. Os valores de H variam de 0 a 1. Valores elevados de H (> 0.80) sugerem uma variável latente bem definida, que provavelmente se mantém estável em diferentes estudos.

\subsubsection*{\textit{Validade baseada na relação com medidas externas}}

Buscou-se avaliar as evidências de validade do instrumento a partir da relação com outras variáveis dispostas em outros questionários semelhantes \cite{Souza2017}, bem como a assimetria de informações. Para isso, foram realizados testes de validade convergente e discriminante. A aplicação da validade convergente visa compreender a relação significativa entre duas ou mais medidas de um mesmo construto, ou construtos relacionados teoricamente \cite{Cunningham2016ImplicitValidity}. A validade discriminante é necessária, pois os construtos para serem teoricamente válidos não devem estar relacionados entre si \cite{Hubley2014DiscriminantValidity}. Nesse sentido, a aplicação de testes de validação descritivos visa avaliar a diferença entre as medidas dos diferentes construtos presentes no instrumento.

\subsubsection*{\textit{Validade convergente e validade discriminante}}

A validade convergente dos fatores foi avaliada utilizando a variância média extraída (AVE) e utilizando cargas fatoriais padronizadas e o erro da medida do item. A estimativa do AVE foi obtida pelas cargas fatoriais e erros de medida obtidos pela análise fatorial confirmatória, utilizando o estimador ponderado da média e variância dos mínimos quadrados (WLSMV). Valores de AVE superiores a 0.50 indicaram boa validade convergente \cite{Maroco2014}. 

\subsubsection*{\textit{Validade da Invariância das Medidas}}

A análise fatorial confirmatória multigrupo (AFCMG) foi realizada para verificar a invariância das medidas do modelo proposto em duas amostras selecionadas aleatoriamente na extensão de 50\% para cada (T=Teste e A=Amostra) dos dados coletados \cite{Borsa2018}. Os testes de Índice de Ajuste Comparativo (CFI) e a diferença do qui-quadrado foram utilizados para avaliar a invariância da escala. Para assumir a invariância da medida, em comparação com o modelo anterior, o índice de ajuste do modelo de teste no CFI não deve ser superior a 0.01 \cite{Cheung2002EvaluatingInvariance}. O ${\chi}^2$ deve apresentar um valor de $p \geq 0.05$ \cite{Milfont2010TestingResearch.,Zait2011MethodsValidity}. 

\subsubsection*{\textit{Validade com base no padrão de resposta do item}}

As propriedades psicométricas do instrumento adaptado foram avaliadas usando o modelo Rasch para dados politômicos \cite{Andrich1978ACategories}. Além disso, foram avaliados indicadores de confiabilidade de pessoas e itens e desvios de desempenho por meio dos índices \textit{Infit} e \textit{Outfit}. Em relação aos índices de confiabilidade, são esperados valores superiores a 0.70 ou, preferencialmente, valores superiores a 0.80 \cite{Linacre2012WinstepsGuide}. Os índices \textit{Infit} e \textit{Outfit} quantificam os resíduos dos itens do modelo testado \cite{Linacre2012WinstepsGuide,Bond2013ApplyingSciences}. A DIF foi avaliada pelo procedimento de Mantel \cite{Wright1994ReasonableTransactions}. Itens cujas estimativas de dificuldade foram significativamente diferentes para homens e mulheres $p \leq 0.05$ foram inspecionados. 
A magnitude do DIF foi interpretada através do contraste DIF: valores entre |0.00| e |0.43| são considerados baixos/insignificantes; valores entre |0.44| e |0.64| são considerados moderados; e valores acima |0.64| são considerados altos \cite{Linacre2012WinstepsGuide}. 

\section{Resultados e discussão}

\subsection{Tradução transcultural dos itens para o instrumento}

O processo de tradução e adaptação cultural gerou a versão em português do Brasil. Durante o processo de tradução, não foram encontradas dificuldades importantes. No entanto, a forma observou alguns problemas idiomáticos, resultando em algumas mudanças fracas nas estruturas gramaticais. Com base no consenso dos tradutores, foram aceitos 23 itens adaptados, 19 politômicos (escala Likert 1 - 5), dois dicotômicos (Sim - Não) e uma questão qualitativa também foi aceita.


\subsection{Validade do conteúdo}

Com base nas pontuações atribuídas aos itens, foi calculado o CVC para avaliar a concordância que os juízes realizam na escala. Com base na pontuação atribuída pelos juízes, os itens que não atingiram valor maior ou igual a 0.80 foram agrupados e realocados para correção em pelo menos dois dos critérios estabelecidos (clareza, relevância e relevância). Ao implementar os procedimentos de validação das normas APA para adaptação transcultural na população de língua portuguesa, o estudo qualitativo inicial, como esperado, revelou sobreposições iniciais de tradução, bem como discrepâncias importantes nos idiomas utilizados. 

A partir do instrumento base em sua versão original, cinco alterações foram aceitas nas afirmações e submetidas a novas rodadas. RI 02 \textit{“Você acredita que esse símbolo/logotipo da marca é fácil de ser lembrado?”} foi corrigido para \textit{“Você acha que este símbolo/logotipo da marca é fácil de lembrar”} (clareza = 0.430); (relevância = 0.650); (relevância = 0.700); RI 01 \textit{“Ao imaginar o produto, percebi que a qualidade dos produtos é alta”} alterada para \textit{“Ao imaginar um produto dessa marca, acredito que a qualidade é alta”} (clareza = 0.896) (relevância = 0.963); (relevância = 0.963); BQP 01 \textit{“Sinto que ao adquirir os serviços, receberei a qualidade que espero”} Alterado para \textit{“Sinto que ao adquirir os serviços desta marca, a qualidade atenderá às minhas expectativas”}, BQP 02 (clareza = 0.830) (relevância = 0.963); (relevância = 0.896) \textit{“Ao visualizar os serviços da Marca, você conseguiria diferenciar os produtos de outras empresas?”} alterado para \textit{“Após ver a imagem acima, você conseguiria diferenciar serviços de marca de outras empresas?”}; PI 03 (clareza = 0.896) (relevância = 0.963); (relevância = 0.963) \textit{“Se os preços da Marca subirem um pouco, eu compraria da concorrência, independentemente da responsabilidade da sustentabilidade ambiental”} Alterado para \textit{“Se os preços da marca subissem um pouco, eu compraria da concorrência, mesmo sabendo que eles se preocupam com a sustentabilidade ambiental”}. Os demais itens listados para compor a escala atenderam aos critérios estabelecidos.

A alteração do item RI 01 foi sugerida considerando a mudança do termo “perceber” para “acreditar” por imaginar que o item possui algum fator de qualidade e está semanticamente mais próximo da crença do que da percepção. Tendo em vista que o CVC apresentou resultados satisfatórios para todos os aspectos julgados, a mudança pareceu adequada, sugerindo que o termo era semanticamente adaptável.

Após as correções, o Coeficiente de Validade de Conteúdo da escala completa foi aceito para elegibilidade do instrumento para a nova fase de validação. Considerando a média do $CVC_{t}$ para três aspectos julgados — clareza, coerência e relevância - foi $CVC_{t}$ = 0.87. Com base nos valores encontrados, o $CVC_{t}$ para todo o instrumento apresentou valores satisfatórios de clareza de linguagem ($CVC_{t}$ = 0.84), relevância prática ($CVC_{t}$ = 0.91) e relevância teórica ($CVC_{t}$ = 0.90), tendo um coeficiente de validade de conteúdo para o teste na totalidade de $CVC_{t}$ = 0.87 e um erros P= 0.003. O instrumento foi considerado elegível para as demais fases, considerando os valores semelhantes aos encontrados por  \citeonline{Daffern2017PredictorsPunctuation:}.

\subsection{Validade com base na estrutura interna}

Valores de Bartlett (2926.2, gl = 210; p < 0.001) e KMO (0.890) sugeriram a interpretabilidade da matriz de correlação dos itens \cite{Hutcheson1999TheScientist,Hutcheson1999TheScientist}. Para estabilidade, o alfa de Cronbach apresentou valores aceitáveis (Cronbach's = 0.925). A análise paralela sugeriu mais quatro fatores representativos para os dados Tabela \ref{tab:5}.

\begin{table}[H]
\centering
\caption{Resultados da análise paralela da escala de percepção dos consumidores sobre marcas sustentáveis na Amazônia}
\label{tab:3}
\begin{tabular}{cc}
\hline
\textbf{\begin{tabular}[c]{@{}c@{}}Porcentagem de explicação\\ de dados reais\end{tabular}} &
  \textbf{\begin{tabular}[c]{@{}c@{}}Porcentagem de explicação\\ de dados aleatórios (IC 95\%)\end{tabular}} \\ \hline
49.527* & 10.780 \\
8.888*  & 9.832  \\
7.545*  & 9.026  \\
5.293*  & 8.476  \\ \hline
\end{tabular}
\begin{threeparttable}
\begin{tablenotes}
- [*] Nota: O número de fatores a serem retidos é quatro, pois quatro fatores apresentam  \% de dados reais mais alta do que os dados aleatórios.
\end{tablenotes}
\end{threeparttable}
\end{table}

As cargas fatoriais dos itens podem ser observadas no Apêndice E.  Além disso, são relatados os índices de confiabilidade composta, bem como as estimativas de replicabilidade dos escores fatoriais (H - index) \cite{Ferrando2018AssessingAnalysis}. A partir da análise fatorial, foi possível identificar quais deles pertencem a qual dimensão, ou seja, reorganizá-los contando os itens que melhor os explicavam, bem como quais itens não atingiram cargas fatoriais para serem mantidos no instrumento.

A estrutura fatorial rendeu poucos carregamentos cruzados. Houve apenas um item (\textit{Branding Inovation}): “A marca oferece produtos mais inovadores do que outras empresas que conheço”, que não projetou um valor justo no fator esperado (ou seja, itens com cargas fatoriais acima de 0.30). Baseado em \citeonline{Nunnally1978AnMeasurement}, os itens com fator de carga abaixo de 0.30 foram omitidos. Nos demais, apenas um item apresentou padrão de cargas cruzadas, sendo este item alocado à classe com maior valor de carga e corrobora os achados da literatura. Todos os itens apresentaram valores próximos à carga da própria dimensão, corroborando também a validade discriminante, comportamento semelhante ao encontrado por \citeonline{Castro2020CognitiveMeasurement}.

A confiabilidade composta dos fatores foi aceitável (acima de 0.70) para o fator SBP e marginalmente aceitável. A replicação foi acatada por meio dos valores de H-index <0.8 \cite{Ferrando2018AssessingAnalysis}.

Na confiabilidade composta, os fatores BQP, PI e RI apresentaram valores marginalmente aceitos pela literatura. Segundo \citeonline{Valentini2016AverageCoefficients}, a mudança na homogeneidade das cargas fatoriais pode ser decorrente do menor número de itens no fator para limitar a interpretação dos resultados da confiabilidade composta.
Os índices de ajuste da escala foram adequados (${\chi}^2$ = 86.460. gl = 132; \textit{p} < 0.001; RMSEA = 0.027; BCa 95\% [0.00 – 0.037]; CFI = 0.998; BCa 95\% [0.998 – 1.002]; TLI = 0.997; BCa 95\% [0.997 — 1.002]). Na raiz do erro quadrático médio de aproximação (RMSEA), um valor inferior a 0.08 para o instrumento proposto indica um bom ajuste estatístico para o modelo \cite{Tucker1973AAnalysis,West2012ModelModeling}. Índices de ajuste comparativo (CFI) e índice Tucker-Lewis (TLI) acima de 0.95 também indicam um bom ajuste do instrumento.

É fundamental destacar que os indicadores de unidimensionalidade Congruência Unidimensional (UniCo) 0.845, Variância Comum Explicada (ECV) 0.886 e Média das Cargas Residuais Absolutas do Item 0.177 não corroboraram a hipótese de unidimensionalidade da escala, uma vez que os valores de UniCo, ECV e MIRAEL maiores que 0.95 sugerem que os dados podem ser tratados como essencialmente unidimensionais; caso contrário, aceita-se a hipótese nula \cite{Ferrando2017}, pois valor de MIREAL inferior a 0.300 sugere que os dados podem ser tratados como essencialmente unidimensionais. Portanto, isso implica que nosso instrumento proposto atende aos referenciais para a expressão teórica dos construtos formativos.

A Tabela \ref{tab:8} detalha a qualidade e eficácia dos fatores. O índice de determinação de fatores (IDE) pode ser interpretado como o número de diferentes níveis de fatores que podem ser diferenciados com base nas estimativas de pontuação dos fatores. Valores acima de 0.90 são considerados bons para o IDE \cite{DiStefano2018FittingData:,Shi2020TheIndices.}

\begin{table}[H]
\centering
\caption{Qualidade e eficácia das estimativas de pontuação fatorial}
\label{tab:8}
\begin{tabular}{ccccc}
\hline
\textbf{\begin{tabular}[c]{@{}c@{}}Métodos de estimação* \end{tabular}} &
  \textbf{\begin{tabular}[c]{@{}c@{}} SBP\end{tabular}} &
  \textbf{\begin{tabular}[c]{@{}c@{}} PI\end{tabular}} &
  \textbf{\begin{tabular}[c]{@{}c@{}} RI\end{tabular}}&
  \textbf{\begin{tabular}[c]{@{}c@{}} BQP\end{tabular}} \\ \hline
  IDE & 0.974 & 0.951 & 0.950 & 0.948 \\
  ORION & 0.951 & 0.904 & 0.902 & 0.878 \\
  \hline
\end{tabular}
\begin{threeparttable}
\begin{tablenotes}
- [*] Nota: IDE = Índice de Determinação Fatorial; ORION = Confiabilidade geral de pontuações N-EAP Oblíquas Prévias totalmente informativas.
\end{tablenotes}
\end{threeparttable}
\end{table}

Os valores obtidos pela estrutura fatorial exploratória confirmam a previsão inicial proposta nesta pesquisa para a construção do instrumento, estabelecendo o modelo final com quatro fatores. Ao fator percepção da qualidade (BQP), foram atribuídas cinco questões que descrevem a percepção do potencial consumidor entrevistado quando a marca avaliada apresentou o produto. Para o fator de posicionamento estratégico de marca (SBP), foram atribuídas dez questões. A dimensão da SBP baseia-se teoricamente no conceito de posicionamento estratégico descrito por \citeonline{Lopez2000EstrategiaComparado}, em que os autores definem o posicionamento estratégico da marca como sua atuação na busca de fatores que compõe o sucesso do negócio, como inovação tecnológica e relação com os conceitos éticos do consumidor e do meio ambiente. Por fim, o fator intenção de compra - PI foi composto por 2 questões, e o fator inovação no varejo - RI foi composto por 4 questões.

\subsection{Validade baseada em relações com medidas externas}
\subsubsection*{\textit{Validade convergente e discriminante}}

O modelo proposto a partir da validade por medidas internas compreende um total de 22 itens divididos em quatro dimensões, a saber, SBP (0.845), BQP (0.964), PI (0.889) e RI (0.677). De acordo com \citeonline{Maroco2014}, todas as dimensões do instrumento atingiram valores aceitáveis (> 0.50). A análise par a par apresentou bons valores para atestar a validade discriminante dos fatores da escala presente na (Tabela \ref{tab:9}). O valor VEM para cada dimensão a ser aceita deve ser maior que 0.50. Além disso, as medidas operacionais obtidas pelos fatores, quando analisadas pareadas, apresentaram valores satisfatórios para atestar a hipótese de validade ao analisar a validade discriminante do instrumento.

\begin{table}[H]
\centering
\caption{Validade discriminante dos fatores da escala}
\label{tab:9}
\begin{tabular}{cccc}
\hline
\textbf{\begin{tabular}[c]{@{}c@{}}Correlação de dimensões \end{tabular}} &
  \textbf{\begin{tabular}[c]{@{}c@{}} Coeficiente de correlação Z\end{tabular}} & \textbf{\begin{tabular}[c]{@{}c@{}} $R^2$ \end{tabular}} \\ \hline
\ $SBP \leftrightarrow BQP$ \ & 0.699 & 0.488 \\
\ $SBP \leftrightarrow PI$ \ & 0.846 & 0.715 \\
\ $SBP \leftrightarrow  RI$ \ & 0.857 & 0.734 \\
\ $BQP \leftrightarrow PI$ \ & 0.596 & 0.355 \\
\ $BQP \leftrightarrow  RI$ \ & 0.672 & 0.451 \\
\ $BQP \leftrightarrow  RI$ \ & 0.756 & 0.571 \\
\hline
\end{tabular}
\end{table}

\citeonline{Fornell1981EvaluatingError} explicam que, quando um dos dois fatores (\textit{fej}) são maiores ou iguais ao quadrado da correção entre esses fatores, a validade discriminante é atestada. Então, os fatores explicam mais variabilidade dos itens que são manifestações de cada um dos fatores do que a variabilidade que apenas um fator explica \cite{Borsa2018}. Atingindo índices satisfatórios, o instrumento foi considerado elegível para as demais validações.

 \subsubsection*{Invariância da medida}

A Tabela \ref{tab:10} mostra a análise de invariância multigrupo configural, métrica e escalar. Esses são os resultados da escala nos diferentes grupos investigados. Os dados de diferentes testes do Índice de Ajuste Comparativo ($\Delta$CFI) indicam que a medida apresenta ajustes finos de invariância. Para assumir a invariância de medição, o valor atribuído a ($\Delta$ CFI) não deve ser maior que 0.01. Observou-se, por meio dos valores de $\Delta$CFI, que a estrutura da escala é estável, sem viés de resposta nas amostras para os diferentes grupos \cite{Lee2020FitModel}.



\begin{table}[H]
\caption{Validade discriminante dos fatores da escala}
\centering
\label{tab:10}
\resizebox{\columnwidth}{!}{%
\begin{tabular}{ccccccll}\hline  
\backslashbox{Questionários}{Grupos}
& X2 & \textit{df} &\textit{P-Value} & RMSEA(IC90\%) & CFI & TLI & $\Delta$ CFI \\\hline

Meeting Room        & 2413.245 & 421 & 0.00 & 0.079 (0.062 0.086) & 0.952 & 0.968 & --- \\
Invariância métrica & 1177.096 & 383 & 0.00 & 0.072 (0.065 0.079) & 0.951 & 0.997 & 0.001 \\
Invariância escalar & 1186.748 & 400 & 0.00 & 0.065 (0.059 0.072) & 0.953 & 0.926 & -0.001 \\
\hline

\end{tabular}%
}
\end{table}

Os dados de invariância por análise fatorial confirmatória multigrupo (AFCMG), através dos valores de invariância configurável, métrica e escalar, mostram que as estruturas medidas entre os grupos e as unidades escalares são as mesmas para os grupos estudados, uma vez que os valores de diferença do Índice de Ajuste Comparativo ($\Delta$CFI) está abaixo de 0.01. A invariância fatorial é um componente essencial do processo iterativo de demonstrar a equivalência de medição de construtos latentes entre grupos, incluindo subpopulações de gênero e idade.
Quando a equivalência de medição está presente, a relação entre a variável latente e a variável observada permanece inalterada nas populações da análise \cite{Raju2002MeasurementTheory}. Essa premissa demonstra sucintamente a importância da equivalência de medidas em instrumentos psicométricos, sugerindo ser possível aplicar essa escala quando o objetivo é analisar, em diferentes grupos, as percepções dos consumidores sobre as marcas relacionadas à sustentabilidade ambiental na Amazônia.

\subsubsection*{Validade com base no padrão de resposta do item}

A validade baseada no padrão de resposta aos itens utilizando a TRI de itens, quando associada a análises de invariância, permite avaliar a similaridade dos itens em um determinado instrumento desenvolvido para diferentes grupos \cite{Sireci2021UsingTest}. As análises iniciais mostraram índices de confiabilidade adequados para os itens (Confiabilidade = 0.98; Índice de Separação = 7.48) e para as pessoas (Confiabilidade = 0.90; Índice de Separação = 2,93), o que sugere que as estimativas obtidas tendem a ser replicáveis. Em relação aos desvios de desempenho ( \textit{infit} e \textit{outfit}), vale ressaltar que o instrumento apresenta valores aceitáveis, \cite{Bond2013ApplyingSciences,Boone2016RaschHow}. Tanto para os itens ( \textit{Infit} MNSQ = 1.00. ZSTD = -0.70; \textit{Outfit} MNSQ = 1.06; ZSTD = -0.13) quanto para os participantes ( \textit{Infit} MNSQ = 1.10. ZSTD = -0.10; \textit{Outfit} MNSQ = 1.06; ZSTD = -0.20).

Os resultados da TRI corroboram a validade e confiabilidade do instrumento. Além disso, os valores de confiabilidade dos itens (0.98) e o índice de separação dos itens (7.48), assim como o índice de confiabilidade dos participantes = 0.90 e o índice de separação dos participantes (2.93), sugerem que as estimativas obtidas tendem a ser replicáveis. Os valores de \textit{infit} e \textit{outfit} do instrumento dos itens e dos participantes se mostram aceitáveis \cite{Bond2013ApplyingSciences,Boone2016RaschHow}.

O mapa item-pessoa do teste piloto (Figura \ref{fig:20}) mostra que, no instrumento apresentado, o nível do traço latente foi maior que a dificuldade dos itens (média teta = 0.67; média de dificuldade dos itens = 0). Como visto, os itens mais difíceis foram BQP 02: “Após ver a imagem acima, você consegue diferenciar os serviços desta empresa de outras empresas?” (\textit{d} = 1.83) e RI 04: "Você acredita que este símbolo/logotipo da marca é fácil de lembrar?" (\textit{d} = 1.43), e o item mais específico foi o item PI 01, “eu consumiria essa marca” (\textit{d} = -0.96).


\begin{figure}[H]
  \centering
  \caption{Mapa Pessoa-item teste piloto}
  \includegraphics[scale=0.8]{Imagens/fig_20.pdf}
  \fonte{\textbf{Elaborado pelos autores} (2022)}
  \label{fig:20}
\end{figure}


A maioria dos itens está agrupada entre – 1.5 e aproximadamente 1.5 logit (Apêndice F), portanto há maior precisão na estimativa teta dos participantes alocados nesse intervalo de contínuo, e a Curva de Informações do Teste corrobora essa informação. Em relação às estatísticas dos itens individualmente. O Apêndice F mostra os indicadores de limiar quando todos os itens apresentaram uma estrutura crescente, conforme o teoricamente esperado.

Em relação ao DIF para faixas etárias, O Apêndice G apresenta os resultados encontrados. Como visto, apenas o RI 04 \textit{“Você acha que este símbolo/logo da marca é fácil de lembrar?”} apresentou DIF significativo com tamanho de efeito moderado (Contraste DIF = 0.48). Os itens com maior dificuldade (BQP 02 e RI 04) refletem que os potenciais consumidores precisam de um traço latente maior para confirmação positiva. Esse resultado mostra que, para os consumidores, os signos visuais utilizados pelas marcas não refletem o posicionamento estratégico buscado.

Para entender melhor essa dificuldade nos itens BQP 02 e RI 04, os resultados da análise diferencial de funcionamento do item (DIF) sustentam a premissa de que estes itens apresentam valores de DIF variados quando valores de \textit{b} são observados para sujeitos com diferentes faixas etárias \cite{CallegaroBorsa2012Cross-culturalConsiderations}. As perdas de reconhecimento de termo em adultos podem estar ligadas à atividade realizada em situações de atenção dividida \cite{Craik2018AgeMemory:}, fato que pode ter afetado os participantes da faixa etária mais avançada durante a aplicação do questionário em ambiente digital.

A Figura \ref{fig:9} ilustra a relação entre a curva característica do instrumento e os erros padrão dos itens do instrumento. A linha sólida de Informação do Instrumento ($\theta$) representa a função de informação da escala. A escala geral forneceu mais informações no intervalo 0 $\leq$ $\theta$ -2 $\leq$ +4, caindo para valores próximos a 6. A linha pontilhada de erro padrão do instrumento fornece uma referência visual de como a precisão da estimativa varia em ($\theta$), com valores menores correspondendo a melhores precisão da estimativa. \cite{Bean2021ItemDevelopment}. Para erros, a acurácia do escore estimado do instrumento proposto é ajustada no intervalo de traço latente entre 0 $\leq$ $\theta$ $\leq$ +6.

Outro item obrigatório, que pediu um nível mais alto de traço latente (endosso positivo entre concordo e concordo), foi o PI01, em que o consumidor está disposto a pagar um preço maior (\textit{premium}) até o limite de 10\%. Os resultados aproximam-se dos obtidos ao \citeonline{DAmico2016ExploringSulfites} analisar a probabilidade de um consumidor pagar (ou não) um prêmio por um produto com certificação sustentável e que a sustentabilidade do \textit{Brand Equity} individualmente não promove o interesse do consumidor em pagar mais pelo produto, corroborando os achados de \citeonline{Rahmani2019AreImpacts}.


\begin{figure}[H]
  \centering
  \caption{Curva de informação do Instrumento}
  \includegraphics[scale=0.3]{Imagens/fig_9.pdf}
  \fonte{\textbf{Elaborado pelos autores} (2022)}
  \label{fig:9}
\end{figure}

As funções de informação para cada item são mais bem ilustradas pelas curvas características do item (ICCs), conforme mostrado na Figura \ref{fig:10}a. Essas curvas mostram que a informação do item não é uma quantidade estática, mas condicional aos níveis de dificuldade do item. Por exemplo, conforme mostrado no mapa item-pessoa, o item PI 01 mostra um progresso mais significativo em relação ao nível de traço latente do entrevistado ($\theta$) para respostas positivas ao questionário (concordo e concordo).

O estudo contribui para a literatura de \textit{Branding} ao propor e verificar um modelo que analisa as associações de variáveis latentes que formam o \textit{Brand Equity} (percepção de qualidade da marca, posicionamento estratégico da marca, intenção de comprar e inovação no varejo) em marcas que desenvolvem atividades sustentáveis. A partir de um modelo assim proposto, é possível compreender a rede nomológica que conecta os aspectos de \textit{Branding} e sustentabilidade de empreendimentos ecologicamente responsáveis.

\begin{figure}[H]
\caption{Curvas características do item (a) e da escala (b)}
\center
\subfigure[ref1][Curvas características do item]{\includegraphics[scale=0.15]{Imagens/fig_10.pdf}}
\qquad
\subfigure[ref2][Curva característica da escala]{\includegraphics[scale=0.15]{Imagens/fig_11.pdf}}
\label{fig:10}
\end{figure}

A Figura \ref{fig:10}b descreve a curva característica da escala CCE. A curva mostra que, para o instrumento analisado, um escore $\theta$= -1 se traduz em um escore esperado de 30 pontos, enquanto um valor de $\theta$= 3.4 se traduz em um escore válido esperado de 90 pontos.
Por outro lado, o item SBP 07, que apresentou o menor nível de endosso positivo (concordo e concordo), mostra que o consumidor, após conhecer todas as informações sobre a marca e concordar com o valor dos produtos, mais no máximo 10\% de aumento do prêmio, apresenta uma tendência de consumo das marcas expostas na pesquisa. É importante notar que os itens SBP 07 ( $\theta$= -1.28) PI 02 ($\theta$= -1,18) apresentaram um nível baixo ($\theta$) que sinaliza endossos positivos na escala Likert (concordo e concordo totalmente) e mostra que, após enfrentar todas as informações da marca adicionadas ao acordo com os preços dos produtos, esse potencial consumidor lucra um aumento máximo de 10\% no prêmio, demonstrando desejo de consumir as marcas apresentadas na pesquisa.

A curva de informações do instrumento explica brevemente como os itens, em geral, fornecem informações estatísticas sobre o traço latente. As curvas são funções matemáticas umas das outras onde Erro Padrão ($\theta$) = 1 – $\surd(\theta$). Mais informações ao longo da escala ($\theta$) levam a menores erros padrão, resultando em estimativas ($\theta$) mais precisas. Além disso, os valores das informações da escala podem ser usados para calcular qual nível de ($\theta$) o instrumento se ajusta melhor de forma confiável. A amplitude de informação obtida pelos itens validados estando entre 0 $\leq$ $\theta$ $\leq$ +4 indica que o instrumento se aplica melhor a potenciais consumidores que tendem a responder moderadamente (entre discordo e concordo).

A transformação das estimativas de endosso presentes na Curva Característica do Instrumento permite verificar que os participantes do teste, em média, responderam 50\% acima da pontuação positiva da escala $\leq$ 4 (concordo na escala Likert). A transformação para a métrica da escala original fornece um referencial mais familiar para interpretar a soma dos escores, \cite{Bean2021ItemDevelopment}, possuindo propriedades favoráveis à confiabilidade na análise do traço latente, melhorando e tornando a abordagem em escores mais precisa \cite{Edwards2009MeasurementChange}. Em relação aos efeitos do teto e piso, não foram identificados participantes ou questões que corroborassem essa ocorrência para este instrumento.

A escala final com 22 itens, disponível em versão eletrônica, apresenta boas propriedades psicométricas. Questionários-piloto baseados no trabalho de  \citeonline{Netemeyer2004DevelopingEquity}, \citeonline{Lin2015ConceptualizingTaiwan} e \citeonline{Yoo2017AnEquity:} em um total de 32 itens abrangendo cinco dimensões, foram reduzidos para 20 itens abrangendo quatro dimensões. Cada dimensão contém itens sobre a percepção da qualidade do produto, o posicionamento percebido da marca, a intenção do potencial consumidor em comprar e a inovação percebida no \textit{Branding} da marca analisada.

O processo de desenvolvimento de instrumentos psicométricos por adaptação transcultural é uma abordagem que pode ser aplicada a diferentes ambientes culturais e linguísticos \cite{Cozendey-Silva2016Cross-culturalTool}. A inclusão de países de língua portuguesa pode ajudar a preencher a lacuna de dados na rede nomológica que representa a percepção do consumidor sobre o valor da marca sustentável. O objetivo da adaptação é tal que o objetivo inicial deste instrumento é alcançar a equivalência entre os instrumentos das medidas originais e suas versões adaptadas \cite{Beaton2007RecommendationsMeasures}. Neste sentido, a versão proposta para o idioma português deste instrumento foi obtida por cuidadosos passos para validação dos instrumentos descritos nas recomendações para padrões de teste APA \cite{Association2018}.

Além dos dados de validação do instrumento, os resultados iniciais sugerem que, dentre as dimensões das marcas sustentáveis, elas exercem influência positiva na variável resultado do valor da marca para o consumidor quando se agrega valor \textit{premium} aos produtos.

A versão em português da proposta tem potencial para gerar indicadores e \textit{feedback}, pois os achados do estudo estão em linha com os resultados de autores anteriores \cite{Ng2014AntecedentsApproach,Shanti2021}, sugerindo que esta escala consegue avaliar dados semelhantes aos propostos em escalas anteriores em inglês. No entanto, uma investigação dos dados resultantes por meio de pesquisas de comparação experimental deve ser realizada futuramente, corrigindo-os, medindo-os e tornando esse instrumento mais robusto \cite{Stavropoulos2019AParents}.

Embora esta pesquisa tenha o objetivo de desenvolver e validar uma escala psicométrica, os presentes achados contribuem para as descobertas científicas que estudam marcas sustentáveis, pois planos de \textit{Branding} eficazes requerem amplo conhecimento sobre a percepção dos consumidores e potenciais consumidores quanto às dimensões responsáveis pelo posicionamento estratégico do patrimônio da marca.

\section{Conclusões}

A escala proposta mostrou-se psicometricamente sólida, apresentando relações esperadas com as variáveis conhecidas quando associada à percepção de potenciais consumidores no tocante à validade e confiabilidade por padrões. Portanto, o instrumento validado pode ser uma valiosa ferramenta de pesquisa na observação da percepção dos consumidores brasileiros em diferentes grupos de sua população e na avaliação da gestão de \textit{Branding} para futuras marcas sustentáveis.

Este estudo produz várias recomendações relacionadas à pesquisa de \textit{Branding} de marca sustentável. Entre eles estão a percepção positiva de produtos sustentáveis, a aceitação de pagar um preço \textit{premium} de até 10\% por produtos de marcas sustentáveis e a diferenciação da memorização de logotipos para diferentes idades.

Essa escala mostrou-se útil para gestores de empresas, que almejam desenvolver seus negócios de forma sustentável e contínua, bem como gestores de marcas voltadas para consumidores de produtos florestais. Estudos futuros podem adotar medidas de investigação por evidências de validade baseadas em teoria do conjunto, incluindo também medidas adicionais de funcionamento cognitivo, modelagem por equações estruturais e Análise Qualitativa comparativa (\textit{Qualitative Comparative Analysis} - QCA) com dados qualitativos dos \textit{stakeholders} envolvidos nas marcas, aplicando-os em exames de grupo para testar ainda mais sua qualidade e utilidade.

Outra recomendação é a de que a autenticidade deva ser incluída em adaptações futuras da escala. E, desta forma, aferir a lealdade da marca e suas influências na predileção e aquisição dos produtos.

Portanto, esperamos que este estudo sirva como ponto de partida para futuras pesquisas para construir e expandir nosso conhecimento acerca desse tema cientificamente interessante e relevante para a gestão.